% !TEX root = ../main.tex
\part{Comando e Comunicação}

Um exército sem comando é apenas uma multidão armada. Esta parte explica como os Pontos de Comando fluem de Oficiais para unidades de suporte que não têm liderança própria, e como um simples Operador de Rádio pode ser a diferença entre um pelotão coordenado e um bando de esquadrões isolados.

\section{Estrutura de Comando}

Unidades de Suporte (HMG, Morteiros, Canhões AT) não possuem líderes orgânicos para gerar PC. Para agir com eficiência total, dependem da cadeia de comando.

\subsection{O Tenente Independente (QG de Pelotão)}

O Tenente não deve fazer parte de um Esquadrão de Fuzileiros --- ele forma uma unidade separada de Comando de Pelotão (Tenente + Rádio + Ajudante). Como raramente precisa se ativar para atirar, ele tem seus 2 PC livres para doar a Unidades de Suporte.

Uma unidade de Suporte pode ser ativada gastando 1 PC de um Oficial dentro do Raio de Comando (15cm para Tenentes, 30cm para Capitães) ou de um Sargento de Pelotão anexado à unidade.

\textbf{Benefício}: a unidade age normalmente, podendo mover-se (se a arma permitir) ou atirar com capacidade máxima (Acerto base 4+).

\subsection{O Sargento de Pelotão (Braço Direito)}

Um Sargento Auxiliar barato pode ser incluído para gerar 1 PC extra, restrito a Unidades de Suporte --- útil quando o Tenente já está ocupado com outras armas.

\subsection{A Regra do Cabo Especialista (Tiro Autônomo)}

Se o jogador não quiser ou não puder alocar um PC a uma unidade de Suporte, ela ainda pode agir de forma autônoma:

A unidade NÃO pode realizar ações de Movimento ou Reposicionamento --- só pode declarar a ordem de ``Fogo''.

\textbf{Penalidade}: -1 no Teste de Acerto por falta de coordenação e observação de um líder.

\begin{nota}
Um único Tenente solto consegue comandar com eficiência total 2 Unidades de Suporte por turno (ex.: 1 HMG e 1 Morteiro), gastando seus 2 PC. Para mais de 2 unidades de suporte sem sargento próprio, promova o Tenente a Capitão ou adicione um Sargento de Pelotão.
\end{nota}

\subsection{O Esquadrão de QG (Comando de Pelotão)}
\label{sec:qg-pelotao}

\textbf{Composição}: 1 Oficial (Tenente ou Capitão) mais até 2 auxiliares --- tipicamente 1 Operador de Rádio e 1 Ajudante/Observador. O grupo todo é tratado como uma única unidade de 2 a 3 miniaturas, movendo-se e mantendo Coesão como qualquer Esquadrão (Parte IV).

\textbf{Geração de PC}: o Oficial gera PC de acordo com seu posto (Tenente = 2, Capitão = 3), independentemente de quantos auxiliares estejam vivos.

\textbf{Movimento}: move-se como Infantaria padrão (15cm / 30cm de corrida).

\textbf{Ativação para Disparo}: o QG raramente precisa se ativar para atirar, já que normalmente gasta seus PC ativando outras unidades (veja as subseções anteriores desta Parte). Se decidir atirar, cada auxiliar armado (pistola, 1 dado) participa normalmente, seguindo a Sequência de Disparo padrão (Parte III).

\textbf{Proteção --- Escudo de Estado-Maior}: ao ser declarado como alvo de tiro, se houver qualquer outra unidade amiga a até 5cm do QG e mais exposta a quem atira, o defensor pode exigir que o inimigo mire nessa unidade mais próxima em vez do QG. Exceção: um Sniper usando Tiro Selecionado (14.6) pode escolher o Oficial diretamente, ignorando esta proteção.

\textbf{Baixa do Ajudante/Observador}: aplique a Regra Geral de Perda do Assistente (\ref{sec:perda-assistente}) --- o Oficial sofre -1 permanente no Teste de Acerto pessoal (relevante apenas se o QG decidir atirar) e qualquer bônus daquele auxiliar deixa de funcionar.

\textbf{Baixa do Operador de Rádio}: a Rede de Comando cai imediatamente --- veja 21.3.

\textbf{Baixa do Oficial}: o Pelotão perde imediatamente os PC daquele posto. Se houver um Sargento de Pelotão (20.2) em campo, ele assume a ativação de Unidades de Suporte normalmente (com sua penalidade padrão). Caso contrário, as unidades de Suporte sem sargento próprio só podem atirar de forma Autônoma (20.3) até a chegada de outro Oficial.

\begin{nota}
Na prática, com o time que você já está jogando (Tenente/Capitão + Observador + Rádio), o Observador segue a regra geral de Perda do Assistente, e o Rádio é o alvo prioritário do inimigo --- proteja-o com o Escudo de Estado-Maior acima.
\end{nota}

\section{Comunicação por Rádio}

Para coordenar ataques em mapas grandes ou manter o oficial móvel, os exércitos utilizam Operadores de Rádio para estender o alcance de suas ordens.

\subsection{Estabelecendo a Conexão}

\textbf{O Transmissor (QG)}: a unidade de Comando (Tenente ou Capitão) deve incluir 1 Operador de Rádio em sua formação.

\textbf{O Receptor (Suporte)}: o Time de Suporte (HMG, Morteiro, Canhão AT) deve possuir a melhoria ``Equipamento de Rádio''.

\subsection{Sinal Aberto (O Benefício)}

Sempre que um Oficial com Operador de Rádio vivo gastar 1 PC para ativar uma unidade de Suporte, ignore a regra de Raio de Comando: o Oficial pode emitir a ordem de qualquer lugar da mesa, independentemente de distância ou linha de visão.

\subsection{Quebrando a Linha (O Ponto Fraco)}

O rádio é uma vantagem tática enorme, o que torna o Operador um alvo prioritário. Se o Operador de Rádio for morto (por exemplo, por um Sniper), a Rede de Comando cai: as unidades de Suporte que dependiam do rádio voltam à regra padrão (o Oficial precisa estar a 15cm, ou elas atiram de forma Autônoma com -1 para Acertar).
